\begin{abstract}
Repeated queries over shared documents repeatedly execute the same lower
transformer layers. We introduce \textbf{CoMem}, which makes split depth $j$ an
explicit reusable-context axis: write one depth-$j$ residual per token, select a
bounded chunk set, and resume only layers $[j{:}L)$. Among document-reuse
systems we are aware of, CoMem jointly makes split depth a tunable serving axis
and isolates it with a matched $j{=}0$ endpoint. On Qwen3-8B, $j{=}12$ reduces
selected-pack Read from 931.9 to 664.4\,ms ($1.403\times$), with a 3.12-point
RULER cost (95\% CI $[2.36,3.93]$); a continuous-prefix oracle recovers the full
gap. The resulting depth axis quantifies a quality--latency--storage trade-off:
combined with bounded selection, a store-ready 128k online-prefill point is
$64.9\times$ faster than dense prefill, while a separate same-adapter,
Write-inclusive pipeline is $2.74\times$ faster. Equal-latency raw replay leads
by 11.56 points with BM25, directly measuring the quality boundary of prepaid
depth rather than hiding it. CoMem stores 8\,KiB/token; a same-Qwen3
CacheBlend-style baseline stores $18\times$ more yet reaches 74.70 versus 97.05
RULER at its strongest tested recomputation ratio. A context--position
factorization identifies missing lower-layer document context as the dominant
tested multikey error, and a 32-token overlap raises 92.5 to 98.5 without
increasing persistent bytes or per-query Read. CoMem opens transformer depth as
a measurable, tunable dimension for repeated-query long-context serving.
\end{abstract}
